\chapter{Experimental Setup}
\label{ch:setup}

\section{Data, models and configurations}

\paragraph{Data.} CIFAR-100~\cite{krizhevsky2009cifar} has 50,000 training and 10,000 test images of
$32\times32$ pixels in 100 classes. We hold out 5,000 training images for validation, which leaves 45,000 for
training. Which images are held out depends on the seed; the test images are the same for every run. The images are
upsampled to $224\times224$ with bilinear interpolation and normalised with the statistics each model was pretrained
with. No data augmentation is used.

\paragraph{Models.} The main experiments use ViT-B/16 (\texttt{google/vit-base-patch16-224}) and Swin-B
(\texttt{microsoft/swin-base-patch4-window7-224}), both pretrained on ImageNet. The closing study also uses DeiT-B
(\texttt{facebook/deit-base-patch16-224}) and, for memory measurements only, ViT-L
(\texttt{google/vit-large-patch16-224}). The 1000-class classifier is replaced by a new, randomly initialised
100-class classifier, which is always trained.

\paragraph{Full fine-tuning (FT).} All weights are trained. The learning rate is the one with the lowest
validation loss among four rates tried (Section~\ref{sec:lr}): $3\times10^{-5}$ for ViT-B and $10^{-5}$ for Swin-B.

\paragraph{LoRA.} Rank $r=8$, scaling $\alpha=16$ and learning rate $10^{-4}$ unless stated otherwise. Three
\emph{targets} are compared, on all blocks: \emph{Attn} (adapters on the query, key and value projections),
\emph{MLP} (adapters on the two linear layers of the MLP) and \emph{Attn+MLP}. Four \emph{placements} of attention
LoRA on half of the blocks are compared with all blocks: the \emph{first half} (blocks 0--5 of ViT-B, 0--11 of
Swin-B), the \emph{second half}, the \emph{even} blocks (0, 2, \ldots) and the \emph{odd} blocks (1, 3, \ldots).
LoRA trains 0.3--1.6 million parameters against 86--87 million for full fine-tuning (Table~\ref{tab:params} in the
appendix).

\paragraph{Training.} All runs use 32-bit floating point, the AdamW optimiser~\cite{loshchilov2019adamw} with
weight decay 0.01, a constant learning rate, batch size 64 and at most 50 epochs. After every epoch the loss on the
validation images is computed. An epoch counts as an improvement only if its validation loss is more than 0.001
below the best so far, and only then are its weights saved. Training stops after 5 epochs without an improvement,
and the saved weights are tested. We call the epoch whose weights are tested the \emph{best epoch} and the last
epoch the \emph{stopping epoch}.

\paragraph{Seeds.} Every configuration of the main experiments is run with the five seeds 42, 123, 456, 789 and
1024; the learning-rate check and the training runs of the closing study use the first three. The seed fixes the
validation split, the initial values of the new weights and the order of the training batches. The runs are
\emph{seeded}, but two runs with the same seed are not bit-identical on the GPU; repeated ViT-B runs differ by up to
a few tenths of a point (Appendix~\ref{app:impl}).

\paragraph{Hardware and software.} All runs use one partition (44.25~GiB) of an NVIDIA B200 GPU, which NVIDIA's
Multi-Instance GPU mechanism divides into partitions; the other partitions were used by other jobs at the same
time. The software is PyTorch~2.11.0~\cite{paszke2019pytorch} (CUDA~12.8, Python~3.12.3) with the HuggingFace
Transformers~5.16.1 and PEFT~0.20.0~\cite{mangrulkar2022peft} libraries, the same versions for every run. Every ViT-B LoRA configuration of Chapters~\ref{ch:lora_ft} and~\ref{ch:placement} was also trained with our own
implementation of the model and of LoRA, from identical initial weights; the two implementations agree closely
(Appendix~\ref{app:impl}), and the tables give the library results. The code, the result files of every run and
the scripts that compute every number in this report are available at
\url{https://github.com/Photon-139/rnd-project}.

\section{Metrics}
\label{sec:metrics}

\begin{itemize}
  \item \textbf{Test accuracy}: accuracy on the 10,000 test images of the weights kept by early stopping.
        Differences between accuracies are given in percentage points (pp).
  \item \textbf{Validation loss}: the cross-entropy loss (lower is better) of the kept weights on the 5,000
        validation images. It decides which weights are tested, and we use it to choose between learning rates, so
        that the test set plays no part in any choice.
  \item \textbf{Mean $\pm$ standard deviation}: all tables give the mean and the sample standard deviation over
        the seeds.
  \item \textbf{Paired difference}: the difference between two configurations computed seed by seed, then
        averaged. It removes the variation that comes from the seed.
  \item \textbf{How differences are described}: we give each paired difference as mean $\pm$ standard
        deviation, with the number of seeds in which it has the same sign, and apply no significance tests. A
        difference whose mean is smaller than its standard deviation is described as \emph{within the spread over
        seeds}; this does not mean that the two configurations are equivalent.
  \item \textbf{Epoch time}: the time of one pass over the 45,000 training images (\emph{train}) or, without
        gradients, over the 5,000 validation images (\emph{validation}), averaged over the epochs of a run. Methods
        are compared by epoch time. The total time of a run depends on the number of epochs before early stopping,
        which varies for reasons unrelated to the cost of an epoch; totals are listed in the appendix and are not
        used for comparisons.
  \item \textbf{Forward, backward and optimiser time}: time per training batch of the forward pass, of the loss
        and backward pass, and of the optimiser step, averaged over all batches of a run.
  \item \textbf{Peak training memory}: the largest amount of GPU memory occupied by tensors at any moment of
        training, as counted by PyTorch, in GiB ($1024^3$ bytes), at batch size 64 unless stated otherwise. It
        excludes the memory of the GPU runtime and memory that PyTorch has reserved but is not using (3--8\% more).
        \textbf{Static memory} is its batch-independent part (weights, gradients, optimiser state) and
        \textbf{activation memory} the part kept for the backward pass (Section~\ref{sec:bg_memory}).
  \item \textbf{Inference memory}: the peak memory occupied by tensors during a forward pass without gradients.
  \item \textbf{Stored size}: what must be stored per task, at 4 bytes per trained value.
  \item \textbf{Unmerged, merged, dense}: a LoRA model with the adapters as separate layers; with the adapters
        added into the weights; a model without adapters.
  \item \textbf{Latency}: the time to classify one test image \emph{end to end} (preparing the image, copying it
        to the GPU, the forward pass, copying the class back); p50 is the median over the test images.
        \textbf{Throughput}: test images classified per second in batches.
\end{itemize}

\section{Measurement procedures}
\label{sec:procedures}

\paragraph{Training runs} record accuracy, validation loss, times and peak and static memory on CIFAR-100.

\paragraph{Memory profiler.} A separate script measures the memory of one training step or one forward pass at
batch sizes from 1 to 256, on random inputs of the shape of the images. These models perform the same operations
whatever the values, so the values do not change the memory: the profiler's peak at batch size 64 agrees with the
training runs to within 0.004~GiB (4~MiB) for every configuration, the peak is the same with a trained adapter as with a
random one, and it is identical across the seeds of a configuration. One profiler run per configuration is
therefore sufficient. The agreement is exact for the ViT-B LoRA configurations and within 0.3~MiB for ViT-B full
fine-tuning; on Swin-B the two differ by up to 4~MiB (0.03\% of the peak) in either direction. The two are not the
same procedure (a training run records the largest value over all batches of an epoch, the profiler the median of
five separate steps), and we did not trace the remaining Swin-B difference.

\paragraph{Merging and latency.} These use the adapter trained with seed 42. For merging we compare the
predictions on all 10,000 test images before and after merging, and after 100 cycles of merging and unmerging. For
latency, the test images are classified one at a time in evaluation mode without gradients; 50 warm-up images are
not timed, 9,950 are, and every timed pass ends by waiting until the GPU has finished. Throughput is measured on
2,048 prepared test images in batches of 1 to 128. Each configuration is measured in three separate processes, two
of which time the unmerged model first, and we report their mean. Full fine-tuning is timed with the pretrained
weights and a new classifier, which have the same shapes, and therefore the same speed, as a fine-tuned model. The
CPU latency uses the same models with 4 threads, 190 timed images and one process.
